\section{From spherical Laplace discrepancy to reflected defect}\label{sec:radial}

We now identify the observable generated by the sample-variance quadratic form. Let $r=n-1$, let $P=I_n-n^{-1}\mathbf1\mathbf1^{\mathsf T}$, and choose $B\in\R^{r\times n}$ with $B^{\mathsf T}B=P$ and $BB^{\mathsf T}=I_r$. If $V$ is uniform on $S^{r-1}$, then every column $b_i$ of $B$ has norm
\[
q:=\|b_i\|=\sqrt{\frac{n-1}{n}},
\]
and $\langle b_i,V\rangle$ has the same distribution as $qV_1$. Put $a=\tau q$ and define
\[
D_X(s):=\frac12R_X(s)
 =\frac12\log\frac{M_X(s)M_X(-s)}{e^{s^2}}.
\]
Under dominance, $D_X\ge0$ on $[-\tau,\tau]$.

The exact definition of the discrepancy in Section~\ref{sec:main-results} gives, for $|t|\le\tau$,
\[
e^{-t^2/2}\E\Psi_r(t\sqrt{Q_n})
 =\E_V\exp\left(\sum_{i=1}^n\left[\log M_X(t\langle b_i,V\rangle)-\frac{t^2\langle b_i,V\rangle^2}{2}\right]\right).
\]
Jensen's inequality, exchangeability, and the symmetry of $V_1$ therefore imply
\begin{equation}\label{eq:integrated-D}
n\E D_X(aV_1)\le\log\bigl(1+\mathfrak D_{n-1,\tau}(Q_n)\bigr).
\end{equation}
This is the precise point at which the spherical observable enters the one-dimensional reflection theory.

For completeness, we record the interior inversion of \eqref{eq:integrated-D}. The density of $V_1$ is
\[
f_{n}(v)=
\frac{\Gamma((n-1)/2)}{\sqrt\pi\,\Gamma((n-2)/2)}
(1-v^2)^{(n-4)/2},
\qquad -1<v<1.
\]
Fix $0<\rho'<a$ and choose $0<h<a-\rho'$. On $[-(\rho'+h),\rho'+h]$ the scaled density of $aV_1$ is bounded below by
\[
m_h:=\frac1a\frac{\Gamma((n-1)/2)}{\sqrt\pi\,\Gamma((n-2)/2)}
\left(1-\left(\frac{\rho'+h}{a}\right)^2\right)^{\max(n-4,0)/2}.
\]
The exponential envelope implies that $D_X$ is Lipschitz there. One may take
\[
L_h:=\frac{K}{e\tau}+\rho'+h,
\]
because $M_X(s)\ge1$ and $|x|e^{\tau|x|}\le(e\tau)^{-1}e^{2\tau|x|}$. If $M_*:=\sup_{|s|\le\rho'}D_X(s)$, two disjoint Lipschitz tents centered at the points where the maximum is attained give
\[
2m_h I_{L_h,h}(M_*)
 \le \frac1n\log\bigl(1+\mathfrak D_{n-1,\tau}(Q_n)\bigr).
 \]
Here the tent area is the nonnegative piecewise function
\[
I_{L,h}(M):=
\begin{cases}
M^2/L,& M\le Lh,\\
2hM-Lh^2,& M>Lh.
\end{cases}
\]
Writing $G(y;L,h)=\sqrt{Ly}$ when $y\le Lh^2$ and $G(y;L,h)=y/(2h)+Lh/2$ otherwise, we obtain
\begin{equation}\label{eq:pointwise-D}
\sup_{|s|\le\rho'}D_X(s)
\le\inf_{0<h<a-\rho'}G\left(\frac{\log(1+\mathfrak D_{n-1,\tau}(Q_n))}{2nm_h};L_h,h\right).
\end{equation}
In particular, for fixed $n,\tau,K,\rho'$, the right-hand side is $O(\sqrt{\mathfrak D_{n-1,\tau}(Q_n)})$ as the discrepancy tends to zero. The strict interior restriction is essential when $n\ge5$, because the spherical density vanishes at the endpoints.
